\problemname{Knapsack in a Globalized World}

Globalization stops at nothing, not even at the good old honest profession of a burglar. Nowadays it is not enough to break in somewhere, take everything you can carry and dart off. No! You have to be competitive, optimize your profit and utilize synergies.

So, the new game rules are:
\begin{itemize}
	\item break only into huge stores, so there is practically endless supply of any kind of items;
	\item your knapsack should be huge;
	\item your knapsack should be full (there should be no empty space left).
\end{itemize}

Damn you, globalization, these rules are not easy to follow! Luckily, you can write a program, which will help you decide whether you should loot a store or not.

\section*{Input}

The input consists of:
\begin{itemize}
\item one line with two integers $n$ ($1 \le n \le 20$) and $k$ ($1 \le k \le 10^{18}$), where $n$ is the number of different item types and $k$ is the size of your knapsack;
\item one line with $n$ integers $g_1, \ldots, g_n$ ($1 \le g_i \le 10^3$ for all $1\leq i\leq n$), where $g_1, \ldots, g_n$ are the sizes of the $n$ item types.
\end{itemize}

\section*{Output}

Output ``{\tt possible}'' if it is possible to fill your knapsack with items from the store (you may assume that there are enough items of any type), otherwise output ``{\tt impossible}''.


